% !TEX program = lualatex
\documentclass[a4paper,11pt]{article}

% !TEX program = lualatex
% Common macros for TTSC Adjunction paper

\usepackage{luatexja}
\usepackage{luatexja-fontspec}
\setmainjfont{HaranoAjiMincho}
\setsansjfont{HaranoAjiGothic}

\usepackage{fontspec}
\usepackage{amsmath,amssymb,amsthm}
\usepackage{mathtools}
\usepackage{tikz-cd}
\usepackage{hyperref}

% Operators
\DeclareMathOperator{\Hom}{Hom}
\DeclareMathOperator{\id}{id}

% Double arrow (avoid undefined \LongRightarrow in some setups)
\providecommand{\LongRightarrow}{\Rightarrow}

% Convenience notation (adjust if C is overloaded)
\newcommand{\HomC}{\Hom_C}

% hyperref-safe PDF bookmarks
\pdfstringdefDisableCommands{%
  \def\Phi{Phi}%
  \def\eta{eta}%
  \def\epsilon{epsilon}%
  \def\Hom{Hom}%
  \def\id{id}%
  \def\mathrm#1{#1}%
  \def\circ{o}%
  \def\Rightarrow{=>}%
  \def\/{/}%
}

% Theorem environments (shared)
\newtheorem{theorem}{定理}
\newtheorem{lemma}{補題}
\newtheorem{definition}{定義}
\newtheorem{remark}{注}



\title{二諦随伴 \texorpdfstring{$F_n \dashv G_{2^n+1}$}{Fn dashv G(2^n+1)} の証明骨格}
\author{千葉 将臣}
\date{2026-02-22}

\begin{document}
\maketitle

\begin{abstract}
二諦（世俗諦／究竟諦）を、擬スライス2-圏 $C/D_{\mathrm{fix}0}$ と固定対象 $D_{\mathrm{fix}0}$ により表現し、
左lax 2-関手 $F_n:(C/D_{\mathrm{fix}0})\to C^D$（爆縮）と右2-関手 $G_{2^n+1}:C\to (C/D_{\mathrm{fix}0})$（奇数すくみ）との擬随伴
$F_n \dashv G_{2^n+1}$ を与える。
さらに、可逆2-射上の離散ノルム $\|\cdot\|:\mathrm{Inv2Cell}\to\mathbb{N}$ による整礎降下条件の下で、誤差セルが有限ステップで恒等へ縮退し、随伴が strict 化する枠組みを提案する。
\end{abstract}

\section{導入}
本稿の目的は、二諦随伴 $F_n \dashv G_{2^n+1}$ を、(i) 擬スライス2-圏、(ii) Debt 2-モナドと EM 2-圏、(iii) lax 2-関手、という部品に分解して定義として完備化し、成立条件を仮定として整理することである。

\section{前提：基底2-圏と固定対象}
\begin{definition}[基底2-圏]
基底2-圏 $C$ は、0-セル・1-射・2-射、および水平／垂直合成と whiskering を備える。
\end{definition}

\begin{definition}[固定対象]
$D_{\mathrm{fix}0}$ を $C$ の固定対象（究竟諦）とする。
\end{definition}

\section{二諦擬スライス2-圏 \texorpdfstring{$C/D_{\mathrm{fix}0}$}{C/Dfix0}}
\begin{definition}[擬スライス2-圏]
擬スライス2-圏 $C/D_{\mathrm{fix}0}$ を以下で定義する。
\begin{enumerate}
  \item 0-セルは $(X,x)$（$x:X\to D_{\mathrm{fix}0}$）。
  \item 1-セル $(h,\alpha_h):(X,x)\to(Y,y)$ は $h:X\to Y$ と可逆2-射 $\alpha_h:y\circ h\Rightarrow x$ からなる。
  \item 2-セル $\Gamma:(h,\alpha_h)\Rightarrow(k,\alpha_k)$ は $\Gamma:h\Rightarrow k$ で、$\alpha_k\circ(y*\Gamma)=\alpha_h$ を満たす。
\end{enumerate}
\end{definition}

\section{Debt 2-モナドと EM 2-圏 \texorpdfstring{$C^D$}{C\string^D}}
\begin{definition}[Debt 2-モナド]
Debt 2-モナド $D$ は 2-自己関手 $D:C\to C$ と 2-自然変換 $\eta:\id_C\Rightarrow D$, $\mu:D\circ D\Rightarrow D$ からなる（公理は標準的2-モナド公理）。
\end{definition}

\begin{definition}[EM 2-圏]
EM 2-圏 $C^D$ の 0-セルは $D$-代数 $(X,\delta:DX\to X)$、1-射・2-射は代数構造と可換なものとする。
\end{definition}

\section{左lax 2-関手 \texorpdfstring{$F_n$}{Fn} と右2-関手 \texorpdfstring{$G_{2^n+1}$}{G(2\string^n+1)}}
\begin{definition}[$F_n$（爆縮）]
左lax 2-関手 $F_n:(C/D_{\mathrm{fix}0})\to C^D$ は、スライスの構造射（空へのベクトル）を段階的に圧縮し、未清算分（Debt）を 0-セル内部状態（$D$-代数構造）へ内在化する。
また、比較2-射 $\mu_{\alpha,\beta}:F_n(\alpha)\circ F_n(\beta)\Rightarrow F_n(\alpha\circ\beta)$ を備え、$\mu_{\alpha,\beta}$ は $C^D$ の代数2-射である。
\end{definition}

\begin{definition}[$G_{2^n+1}$（奇数すくみ）]
右2-関手 $G_{2^n+1}:C\to (C/D_{\mathrm{fix}0})$ は、対象 $Y$ に対し $\mathrm{Stale}_{2^n+1}(Y)$（特定図式の colimit）と普遍射 $\hat y:\mathrm{Stale}_{2^n+1}(Y)\to D_{\mathrm{fix}0}$ を与え、
$G_{2^n+1}(Y):=(\mathrm{Stale}_{2^n+1}(Y),\hat y)$ と定める（構成は仮定として置く）。
\end{definition}

\section{二諦擬随伴 \texorpdfstring{$F_n \dashv G_{2^n+1}$}{Fn dashv G(2^n+1)}}
\begin{theorem}[二諦随伴（擬随伴）]
適切な普遍性（$G_{2^n+1}$）と爆縮性（$F_n$）の仮定の下で、擬自然同値
\[
\Phi_{X,Y}:\Hom_{C/D_{\mathrm{fix}0}}(X,G_{2^n+1}(Y))\simeq \Hom_{C^D}(F_n(X),Y)
\]
が構成でき、擬単位 $\eta$ と擬余単位 $\epsilon$、および修正子 $\lambda,\rho$ により擬三角等式が成り立つ。
\end{theorem}

\begin{remark}
ここでは証明の詳細は、$\Phi$ の順逆構成をそれぞれ普遍性（colimit）と代数構造の内在化により与え、擬自然性と貼り合わせコヒーレンスを確認する形で補う。
\end{remark}

\section{ノルムによる整礎降下と strict 化}
\begin{definition}[離散ノルム]
可逆2-射のクラスにノルム $\|\cdot\|:\mathrm{Inv2Cell}\to\mathbb{N}$ を与える。
特に誤差セル（$\mu,\lambda,\rho$）について、$\|\gamma\|=0\Rightarrow \gamma=\id$ を仮定する。
\end{definition}

\begin{theorem}[strict 化（条件付き）]
反復により誤差セルのノルムが真に降下するなら、有限ステップで $\|\mu\|=\|\lambda\|=\|\rho\|=0$ に到達し、擬随伴は（当該段階で）strict 随伴へ縮退する。
\end{theorem}

\section{今後の課題}
具体構成（$\neg^n$ と $\mathrm{Stale}_{2^n+1}$）を与え、(i) colimit の形状、(ii) anti-collusion 公理の位置づけ、(iii) ノルムの具体例、を明示する。

\end{document}

